\chapter*{Using This Reference}
\addcontentsline{toc}{chapter}{Using This Reference}

Parts I through VII taught Vim's grammar from first principles, deriving each command as an instance of a small number of composable rules rather than presenting commands as a list to memorize. Part VIII inverts that approach deliberately. What follows is the list: every normal-mode command, every Ex command, every motion, operator, text object, register, option, built-in function, regular-expression construct, and command-line argument covered or implied by the preceding chapters, organized as tables for lookup rather than as prose for reading start to finish.

This inversion is intentional rather than a lapse in the book's own stated method. A reference table serves a different purpose than a derivation: it answers "what does this do" or "what was that command called again" in the few seconds such a question deserves, without requiring the reader to re-read the chapter that originally introduced it. Each appendix notes, where relevant, which chapter treats the item in question at greater depth, so that a table entry can serve as an index back into the derivation as easily as it serves as a standalone answer.

Nothing in this part introduces new material. Every command listed here was already covered, directly or as an immediate consequence of the grammar, somewhere in Parts I through VII. A reader who has worked through those parts in order will find Part VIII entirely familiar; a reader using this book primarily as a desk reference, consulting Part VIII directly without having read the derivation first, will find each table sufficient on its own for looking up a specific command's syntax, at the cost of the deeper why that the earlier parts were written to provide.
